% ============================================================
\chapter{Environments and Dependency Management}
\label{ch:environments}
\index{environments}
% ============================================================

\section{Why Isolate Environments?}
\label{sec:why-envs}

A classic problem in ML: ``It works on my machine!'' The causes:
\begin{itemize}
  \item Different Python versions
  \item Incompatible library versions
  \item Missing system dependencies (CUDA, cuDNN)
  \item Unconfigured environment variables
\end{itemize}

\begin{definition}[Virtual environment]
A \textbf{virtual environment}\index{virtual environment} is an isolated
directory containing a specific Python installation and a set of packages,
independent of the system and other projects.
\end{definition}

\begin{center}
\begin{tikzpicture}[
  env/.style={rectangle, draw, rounded corners, minimum width=3cm,
    minimum height=2.5cm, align=center, font=\small},
  pkg/.style={rectangle, draw, fill=blue!10, font=\scriptsize,
    minimum width=2.5cm, rounded corners}
]
  \node[env, fill=green!5] (env1) {};
  \node[above=0.1cm] at (env1.north) {\textbf{Project A}};
  \node[pkg] at ([yshift=0.3cm]env1.center) {scikit-learn 1.3};
  \node[pkg] at ([yshift=-0.3cm]env1.center) {pandas 2.0};
  \node[below=0.1cm, font=\scriptsize] at (env1.south) {Python 3.10};

  \node[env, fill=orange!5, right=2cm of env1] (env2) {};
  \node[above=0.1cm] at (env2.north) {\textbf{Project B}};
  \node[pkg] at ([yshift=0.3cm]env2.center) {scikit-learn 1.5};
  \node[pkg] at ([yshift=-0.3cm]env2.center) {pandas 2.2};
  \node[below=0.1cm, font=\scriptsize] at (env2.south) {Python 3.12};

  \node[below=1.5cm of env1, font=\small] (sys) {System};
  \node[below=1.5cm of env2, font=\small] (sys2) {System};
  \draw[dashed, gray] ([xshift=-1cm]sys.north) -- ([xshift=1cm]sys2.north);
\end{tikzpicture}
\end{center}

\section{venv --- Python's Built-in Tool}
\label{sec:venv}
\index{venv}

\cli{venv} is Python's standard module for creating virtual environments.
It is simple, lightweight, and requires no additional installation.

\begin{codeblock}[title=\textbf{Essential venv commands}]
\begin{minted}{bash}
# Create a virtual environment
python -m venv .venv

# Activate (Linux/macOS)
source .venv/bin/activate

# Activate (Windows)
.venv\Scripts\activate

# Install packages
pip install numpy pandas scikit-learn

# Save dependencies
pip freeze > requirements.txt

# Recreate the environment
pip install -r requirements.txt

# Deactivate
deactivate
\end{minted}
\end{codeblock}

\begin{warning}[title=\textbf{pip freeze pitfalls}]
\cli{pip freeze} includes \textbf{all} transitive dependencies, making the
file fragile. Prefer:
\begin{itemize}
  \item Manually maintaining a \file{requirements.txt} with direct
    dependencies and their versions.
  \item Using \pkg{pip-tools} to generate a locked \file{requirements.txt}
    from a \file{requirements.in}.
\end{itemize}
\end{warning}

\subsection{pip-tools: dependency locking}
\index{pip-tools}

\begin{codeblock}[title=\textbf{pip-tools workflow}]
\begin{minted}{bash}
# Install pip-tools
pip install pip-tools

# requirements.in (direct dependencies only)
# numpy>=1.24
# pandas>=2.0
# scikit-learn>=1.3

# Compile (resolve and lock)
pip-compile requirements.in -o requirements.txt

# Install exactly the locked versions
pip-sync requirements.txt
\end{minted}
\end{codeblock}

\section{Conda --- Environments with System Dependencies}
\label{sec:conda}
\index{Conda}

\textbf{Conda}\index{Conda} manages not only Python packages but also
system libraries (CUDA, MKL, OpenSSL). It is the preferred tool for ML
projects requiring GPUs.

\begin{codeblock}[title=\textbf{Essential Conda commands}]
\begin{minted}{bash}
# Create an environment
conda create -n ml-project python=3.11

# Activate
conda activate ml-project

# Install packages
conda install numpy pandas scikit-learn
conda install -c pytorch pytorch torchvision  # PyTorch channel

# Export the environment
conda env export > environment.yml

# Export without OS-specific dependencies
conda env export --from-history > environment.yml

# Recreate the environment
conda env create -f environment.yml

# List environments
conda env list

# Remove an environment
conda env remove -n ml-project
\end{minted}
\end{codeblock}

\begin{codeblock}[title=\textbf{environment.yml}]
\begin{minted}{yaml}
name: ml-project
channels:
  - pytorch
  - conda-forge
  - defaults
dependencies:
  - python=3.11
  - numpy=1.26
  - pandas=2.1
  - scikit-learn=1.3
  - pytorch=2.1
  - pip:
    - mlflow>=2.8
    - wandb>=0.16
\end{minted}
\end{codeblock}

\begin{bestpractice}[title=\textbf{Conda vs pip}]
\begin{itemize}
  \item Use Conda for system dependencies (CUDA, compilers).
  \item Use pip inside a Conda environment for pure Python packages.
  \item \textbf{Never} mix \cli{conda install} and \cli{pip install} for
    the same package---Conda will not see packages installed by pip and
    vice versa.
\end{itemize}
\end{bestpractice}

\section{Poetry --- Modern Dependency Management}
\label{sec:poetry}
\index{Poetry}

\textbf{Poetry}\index{Poetry} is a modern dependency manager that combines
package management, version locking, and package publishing in a single tool.

\begin{codeblock}[title=\textbf{Poetry workflow}]
\begin{minted}{bash}
# Install Poetry
curl -sSL https://install.python-poetry.org | python3 -

# Create a new project
poetry new ml-project
cd ml-project

# Add dependencies
poetry add numpy pandas scikit-learn
poetry add --group dev pytest black mypy

# Install dependencies
poetry install

# Run in the environment
poetry run python train.py
poetry run pytest tests/

# Export to requirements.txt (compatibility)
poetry export -f requirements.txt -o requirements.txt
\end{minted}
\end{codeblock}

\begin{codeblock}[title=\textbf{pyproject.toml (excerpt)}]
\begin{minted}{toml}
[tool.poetry]
name = "ml-project"
version = "0.1.0"
description = "ML classification project"

[tool.poetry.dependencies]
python = "^3.10"
numpy = "^1.26"
pandas = "^2.1"
scikit-learn = "^1.3"

[tool.poetry.group.dev.dependencies]
pytest = "^7.4"
black = "^23.10"
mypy = "^1.6"
\end{minted}
\end{codeblock}

\section{Tool Comparison}
\label{sec:env-comparison}

\begin{center}
\begin{tabular}{@{}lcccc@{}}
\toprule
\textbf{Criterion} & \textbf{venv+pip} & \textbf{Conda} & \textbf{Poetry} & \textbf{uv} \\
\midrule
Built into Python & Yes & No & No & No \\
System deps & No & Yes & No & No \\
Locking & pip-tools & Yes & Yes & Yes \\
GPU/CUDA & Manual & Yes & Manual & Manual \\
Speed & Fast & Slow & Medium & Very fast \\
Pkg publishing & No & No & Yes & No \\
Learning curve & Low & Medium & Medium & Low \\
\bottomrule
\end{tabular}
\end{center}

\section{uv --- The Fast New Tool}
\label{sec:uv}
\index{uv}

\textbf{uv} is a Python package manager written in Rust, compatible with
pip, extremely fast (10--100$\times$ faster than pip).

\begin{codeblock}[title=\textbf{uv workflow}]
\begin{minted}{bash}
# Install uv
curl -LsSf https://astral.sh/uv/install.sh | sh

# Create environment and install
uv venv
source .venv/bin/activate
uv pip install numpy pandas scikit-learn

# Or all at once
uv pip install -r requirements.txt

# Locking
uv pip compile requirements.in -o requirements.txt
uv pip sync requirements.txt
\end{minted}
\end{codeblock}

\section{Best Practices for Reproducibility}
\label{sec:env-best-practices}

\begin{bestpractice}[title=\textbf{Golden rules for environments}]
\begin{enumerate}
  \item \textbf{One project = one environment}: never share an environment
    between projects.
  \item \textbf{Versioned dependency file}: \file{requirements.txt} or
    \file{environment.yml} in the Git repository.
  \item \textbf{Pin versions}: \cli{numpy==1.26.2}, not \cli{numpy} without
    a version.
  \item \textbf{Separate dev and prod}: development tools (pytest, black)
    should not be in production.
  \item \textbf{Document the Python version}: in the README or
    \file{pyproject.toml}.
  \item \textbf{Never modify the system environment}: never use
    \cli{sudo pip install}.
\end{enumerate}
\end{bestpractice}

\begin{antipattern}[title=\textbf{Common anti-patterns}]
\begin{itemize}
  \item Installing all packages in the system environment
  \item Using \cli{pip install} without \cli{-r requirements.txt}
  \item Forgetting to update \file{requirements.txt} after adding a package
  \item Mixing Conda and pip haphazardly
  \item Sharing a Conda environment via an absolute path
\end{itemize}
\end{antipattern}

\section{Mini-project: Reproducible Environment}
\label{sec:mini-project-ch2}

\begin{miniproject}[title=\textbf{Mini-project 2: Setting up an environment}]
Building on the structured project from Chapter~\ref{ch:introduction}:
\begin{enumerate}
  \item Create a virtual environment with \cli{venv} or Conda.
  \item Write a \file{requirements.in} with direct dependencies.
  \item Generate a locked \file{requirements.txt} with \pkg{pip-tools}.
  \item Write an equivalent \file{environment.yml} for Conda.
  \item Verify that a colleague can recreate the environment from the
    dependency file alone.
  \item Add a \cli{setup.sh} script that automates environment creation.
\end{enumerate}
\end{miniproject}

\section{Exercises}
\label{sec:exercises-ch2}

\begin{exercise}[$\bigstar$ --- Environment creation]
Create a virtual environment with each of the three tools (venv, Conda,
Poetry) for a project using \pkg{numpy}, \pkg{pandas}, and \pkg{matplotlib}.
Compare the environment sizes and creation times.
\end{exercise}

\begin{exercise}[$\bigstar\bigstar$ --- Conflict resolution]
A project requires \pkg{tensorflow>=2.15} and \pkg{numpy<1.24} (a real
conflict). Use \pkg{pip-tools} to identify the conflict, then find
compatible versions. Document the resolution process.
\end{exercise}

\begin{exercise}[$\bigstar\bigstar\bigstar$ --- Cross-platform environment]
Create an ML project that works on Linux, macOS, and Windows:
\begin{enumerate}
  \item Write a multi-OS compatible \file{environment.yml} for Conda
    (\cli{conda env export --from-history})
  \item Create a \cli{setup.sh} / \cli{setup.bat} script that detects
    the OS and installs appropriate dependencies
  \item Test on at least two systems (or use Docker containers to simulate)
  \item Handle the GPU case (install PyTorch CPU vs CUDA depending on the
    machine)
\end{enumerate}
\end{exercise}

\section{Cheatsheet}

\begin{cheatsheet}[title=\textbf{Chapter 2 Summary}]
\begin{tabular}{@{}p{4cm}p{8.5cm}@{}}
\toprule
\textbf{Tool} & \textbf{Key Commands} \\
\midrule
\textbf{venv} & \cli{python -m venv .venv}, \cli{source .venv/bin/activate} \\
\textbf{pip-tools} & \cli{pip-compile requirements.in}, \cli{pip-sync} \\
\textbf{Conda} & \cli{conda create -n env python=3.11}, \cli{conda env export} \\
\textbf{Poetry} & \cli{poetry add pkg}, \cli{poetry install}, \cli{poetry run} \\
\textbf{uv} & \cli{uv pip install}, \cli{uv pip compile}, \cli{uv pip sync} \\
\bottomrule
\end{tabular}
\end{cheatsheet}
